\section{TASA Advantage Properties}
\label{app:tasa_properties}

We provide brief proofs of the four properties stated in \Cref{sec:tasa} for binary rewards $r_i \in \{0, 1\}$ with threshold $c = 0.5$.

\paragraph{Property 1: Sign preservation.}
If $r_i = 1$, then $(r_i - c)_+ = 0.5 > 0$ and $(c - r_i)_+ = 0$, so $A_i = 0.5/\Zp > 0$.
If $r_i = 0$, then $(r_i - c)_+ = 0$ and $(c - r_i)_+ = 0.5 > 0$, so $A_i = -0.5/\Zm < 0$.

\paragraph{Property 2: Monotonicity.}
Since $r_i \in \{0, 1\}$ and $A_i(r_i{=}1) > 0 > A_i(r_i{=}0)$, monotonicity holds trivially.

\paragraph{Property 3: Boundedness.}
The maximum advantage occurs when $r_i = 1$ and all other responses also have $r_j = 1$: $A_i = 0.5/(G \cdot 0.5) = 1/G \leq 1$. The minimum is $-1/G \geq -1$.

\paragraph{Property 4: Non-degeneracy.}
For an all-correct group ($r_j = 1$ for all $j$): $\Zp = G \cdot 0.5$, $\Zm = \epsilon$, and $A_i = 0.5 / (G \cdot 0.5) = 1/G > 0$. For an all-incorrect group: $A_i = -1/G < 0$. No group yields zero advantage.

\section{Implementation Details}
\label{app:impl}

\paragraph{LoRA configuration.}
We apply LoRA to all linear layers (q\_proj, k\_proj, v\_proj, o\_proj, gate\_proj, up\_proj, down\_proj) with rank $r{=}64$, $\alpha{=}128$, dropout 0.05.

\paragraph{Training hyperparameters.}
Learning rate: $2{\times}10^{-5}$; warmup ratio: 0.05; weight decay: 0.01; max gradient norm: 1.0; bfloat16 precision; max completion length: 256 tokens; batch size: 1 per device with gradient accumulation 4.

\paragraph{Replay parameters.}
Replay weight $\lambda = 0.05$; warmup $W = 50$ steps; bank capacity $K = 2$ per prompt; success threshold $r \geq 0.5$ (equivalent to $r = 1$ under binary reward).

\paragraph{Evaluation.}
Greedy decoding with max 256 new tokens. Answer extracted via ``\#\#\#\#'' delimiter regex. Accuracy computed as correct / total evaluated on the full GSM8K test set ($n = 1{,}319$).

\paragraph{Hardware.}
All experiments run on H100 80GB GPUs, one GPU per run. Training takes approximately 3--4 hours per 200-step run.
